\documentclass{article}
\usepackage{notes}

\begin{document}
\stdtitle{PHYSIK II}{Serie 4}{Jakob Schildhauer}

\section{Potentielle Energie einer Ladungsverteilung}
$W = \int dW(r) = \int_{0}^{a} W(r) \, dr = \int_{0}^{a} \frac{4\pi r^2 \rho_0 4 \pi r^3 \rho_0}{12r \pi \epsilon_0} \, dr
=  \int_{0}^{a} \frac{4\pi r^4 \rho_0^2}{3\epsilon_0} \, dr = \frac{4\pi a^5 \rho_0^2}{15 \epsilon_0}$, \\
$Q = \frac{4 \pi a^3 \rho_0}{3} \implies \rho_0 = \frac{3}{4 \pi a^3} \implies W = \dul{\frac{3 Q^2}{20 \pi \epsilon_0 a}}$

\section{Elektrischer Dipol}
\begin{enumerate}[label=\alph*)]
    \item \dots 
    \item \dots 
\end{enumerate}

\section{E-Feld und elektrisches Potential einer Ladungsverteilung}
\begin{enumerate}[label=\alph*)]
    \item $E_{r \leq R} = \frac{4\pi r^3 \rho}{12 \pi \epsilon_0 r^2} = \dul{\frac{r\rho}{3 \epsilon_0}}$, 
    $E_{r \geq R} = \frac{4 \pi R^3 \rho}{12 \pi \epsilon_0 r^2} = \dul{\frac{R^3 \rho}{3 \epsilon_0 r^2}}$
    
    \item \dots 
    
    \item $\Phi_{r \geq R} = \dul{-\frac{R^3 \rho}{3 \epsilon_0 r}}$ (wie bei einer Punktladung), 
    $\Phi_{r \leq R} = - \frac{r^2 \rho}{6 \epsilon_0} + C$, \\
    $\Phi_{r \geq R}(R) = \Phi_{r \leq R}(R) \implies \Phi_{r \leq R} = - \frac{r^2 \rho}{6 \epsilon_0} +  $

    
    \item \dots
    
    \item \dots
\end{enumerate}

\section{Geladener Kreisring}
\begin{enumerate}[label=\alph*)]
    \item \dots 
    \item \dots 
    \item \dots
\end{enumerate}

\section{Wasserstoffatom}
\begin{enumerate}[label=\alph*)]
    \item \dots 
    \item \dots 
    \item \dots
\end{enumerate}
\end{document}